\documentclass[11pt,a4paper]{ctexart}
\usepackage[margin=23mm,top=25mm,bottom=24mm]{geometry}
\usepackage{amsmath,amssymb,amsthm,mathtools,booktabs,tabularx}
\usepackage{enumitem,xcolor,hyperref,listings}
\usepackage{fancyhdr,lastpage}
\hypersetup{colorlinks=true,urlcolor=blue,linkcolor=blue}
\definecolor{dblue}{HTML}{17374D}
\definecolor{bg}{HTML}{F3F6F8}
\lstdefinestyle{python}{language=Python,basicstyle=\ttfamily\scriptsize,
  backgroundcolor=\color{bg},frame=single,breaklines=true,showstringspaces=false,
  columns=fullflexible,keepspaces=true,tabsize=4,rulecolor=\color{gray}}
\lstset{style=python}
\setlength{\parindent}{2em}
\setlength{\parskip}{0.25em}
\newtheorem{theorem}{定理}[section]
\newtheorem{lemma}[theorem]{引理}
\newtheorem{corollary}[theorem]{推论}
\newtheorem{proposition}[theorem]{命题}
\theoremstyle{definition}
\newtheorem{remark}[theorem]{说明}
\newenvironment{proofcn}{\par\noindent\textbf{证明.} }{\hfill$\square$\par}
\pagestyle{fancy}\fancyhf{}
\fancyhead[L]{\small D2：异质 $K=5$ 完整二次证书的有限覆盖}
\fancyhead[R]{\small 主要成果与后续压缩}
\fancyfoot[C]{\small 第 \thepage\ 页 / 共 \pageref{LastPage} 页}
\newcommand{\calP}{\mathcal P}
\newcommand{\calD}{\mathcal D}
\newcommand{\calV}{\mathcal V}
\newcommand{\calM}{\mathcal M}
\newcommand{\one}{\mathbf1}
\newcommand{\Risk}{\mathsf R}
\newcommand{\Obj}{\operatorname{Obj}}
\begin{document}
\begin{center}
{\Large\bfseries D2：异质五分类器稳健投票的有限证书全集}\\[0.35em]
{\large\bfseries 零截距的精确求解、正截距极点覆盖与压缩判据}\\[0.7em]
2026 年 10 月 13 日\quad |\quad 数学研究备忘录\\
\end{center}
\noindent\textcolor{dblue}{\rule{\textwidth}{1.1pt}}
\noindent\textbf{范围及证明状态。}本文仅保留本阶段的主要、独立结果：零截距证书的精确值；任意实系数正截距证书的可构造有限覆盖；及将完整性压缩为有限支配性测试的方法。另外证明了 $K=5$ 的图证书可用一个无 big-$M$ 的小规模混合整数线性规划精确求解。\textbf{不主张}已得到包含几十或几百项的简洁完整闭式；此前的 $U_{98},U_{128},U_{188}$ 只是合法上界，均不得未经证明当成完整公式。

\section{问题与二次对偶}
错误模式为 $x\in\{0,1\}^5$，边际错误率 $q_i\in(0,1/2)$，成对共同错误率的上界 $b_{ij}\ge0$。允许的联合分布为
\[
\calP(q,b)=\left\{\pi\ge0:\sum_x\pi_x=1,\ \sum_x\pi_xx_i=q_i,\ \sum_x\pi_xx_ix_j\le b_{ij}\right\};
\]
假定 $\calP(q,b)\ne\varnothing$，并使用确定性 self-dual 规则 $F(x)+F(\one-x)=1$。其 minimax 错误率记为 $\Risk^*_{\rm det}$。

固定 $F$ 的最坏分布线性规划的对偶为
\begin{align}\label{eq:dual}
\Risk(F)=\min_{z\in\calD_F}L_z(q,b),\quad
L_z(q,b)&=a+\sum_iq_i\beta_i+\sum_{i<j}b_{ij}\gamma_{ij},\\
\calD_F&=\left\{z=(a,\beta,\gamma):\gamma\ge0,\ 
Q_z(x)\ge F(x)\quad\forall x\in\{0,1\}^5\right\},\notag\\[-0.3em]
Q_z(x)&=a+\sum_i\beta_ix_i+\sum_{i<j}\gamma_{ij}x_ix_j.\notag
\end{align}
强对偶成立且最优值可取到。消去 $F$ 得
\begin{equation}\label{eq:Qfeas}
Q_z(x)\ge0,\qquad
\max\{Q_z(x),Q_z(\one-x)\}\ge1\quad(\forall x),\qquad\gamma\ge0.
\end{equation}
对任意满足式 \eqref{eq:Qfeas} 的 $Q_z$，以 $Q_z(x)-Q_z(\one-x)$ 的符号选择失败侧，就可产生被 $Q_z$ 逐点上界的 self-dual 线性阈值规则，因为
\[
Q_z(x)-Q_z(\one-x)
=\sum_i\left(\beta_i+\frac12\sum_{j\ne i}\gamma_{ij}\right)(2x_i-1).
\]
平局可以用一般位置的小扰动打破。因此最优确定性规则可在有限的有符号线性阈值规则集合 $\mathcal W_5$ 中选取。现有精确整数规则枚举确认 $|\mathcal W_5|=1882$，按模型置换为 60 种类型；下述有限性证明实际上不依赖具体的 1882 计数。

\section{新主定理：零截距完全解决，正截距有限覆盖}
\subsection{零截距证书的精确值}
令非负边权 $y=(y_{ij})$ 的图划分函数为
\[
 g(y)=\min_{S\subseteq[5]}
 \max\left\{\sum_{\{i,j\}\subseteq S}y_{ij},
 \sum_{\{i,j\}\subseteq S^c}y_{ij}\right\},
\qquad
B_{\rm graph}(b)=\inf_{y\ge0:g(y)>0}\frac{\sum_{i<j}b_{ij}y_{ij}}{g(y)}.
\]
\begin{theorem}[零截距面的最优值]\label{thm:zero}
在任意可行边际与成对上界下，所有满足式 \eqref{eq:Qfeas} 且 $Q(0)=a=0$ 的二次证书中，最小对偶目标值恰为
\begin{equation}\label{eq:zero}
\boxed{\inf_{Q:\,a=0} L_Q(q,b)=\min\{q_{\min},B_{\rm graph}(b)\}.}
\end{equation}
无须假定 $\sum_iq_i\le1$，也无须假定系数有理或为整数。
\end{theorem}
\begin{proofcn}
当 $a=0$，由 $Q(e_i)\ge0$ 得 $\beta_i\ge0$。设 $s=\sum_i\beta_i$。若 $s\ge1$，有 $L_Q\ge q_{\min}s\ge q_{\min}$。若 $s<1$，互补模式两侧的一阶项均不超过 $s$，所以对所有 $S$ 都有
\[
\max\{\gamma(E(S)),\gamma(E(S^c))\}\ge1-s,
\quad\Longrightarrow\quad g(\gamma)\ge1-s>0.
\]
因此 $L_Q\ge q_{\min}s+B_{\rm graph}(1-s)\ge\min\{q_{\min},B_{\rm graph}\}$。
反向由两类可行证书达到下确界：独裁证书 $Q(x)=x_i$；以及任何 $g(y)>0$ 所对应的 $Q_y(x)=\sum_{i<j}y_{ij}x_ix_j/g(y)$。后者的互补覆盖性正是 $g(y)$ 的定义。\end{proofcn}

\subsection{所有正截距实系数证书的显式有理全集}
定义 $16$ 维系数向量 $z=(a,\beta_1,\dots,\beta_5,\gamma_{12},\dots,\gamma_{45})$，令 $A\in\{0,1\}^{42\times16}$：前 $32$ 行为
\[
A_x=(1,x_1,\ldots,x_5,x_1x_2,\ldots,x_4x_5),
\]
后 $10$ 行为选取各 $\gamma_{ij}$ 的单位行。对固定规则 $F$，右端 $h_F$ 是 $32$ 个 $F(x)$ 与后 $10$ 个零的拼接，因此 $\calD_F=\{z:Az\ge h_F\}$。

\begin{theorem}[参数无关的有限有理证书全集]\label{thm:atlas}
对 $F\in\mathcal W_5$，定义
\begin{equation}\label{eq:vertices}
\calV_F=\left\{
 A_I^{-1}(h_F)_I:\ |I|=16,\ \det A_I\ne0,\ A A_I^{-1}(h_F)_I\ge h_F
\right\},
\qquad\calV_5=\bigcup_{F\in\mathcal W_5}\calV_F.
\end{equation}
则 $\calV_5$ 有限且由有理系数向量构成；对任意满足式 \eqref{eq:Qfeas} 的\emph{实系数}证书 $Q$，存在 $F\in\mathcal W_5$ 及不超过 $17$ 个 $V_k\in\calV_F$、凸组合系数 $\lambda_k$ 和一个衰退方向 $D$，使
\[
Q=\sum_{k=1}^m\lambda_kV_k+D,\quad
m\le17,\quad\lambda_k\ge0,\quad\sum_k\lambda_k=1,
\]
且对任意可行 $(q,b)$ 都有 $L_D(q,b)\ge0$。特别地，
\begin{equation}\label{eq:atlas}
\boxed{\Risk^*_{\rm det}(q,b)=\min_{V\in\calV_5}L_V(q,b).}
\end{equation}
\end{theorem}
\begin{proofcn}
由式 \eqref{eq:Qfeas} 的互补比较，给定 $Q$ 可选择被其上界的 $F\in\mathcal W_5$，所以 $Q\in\calD_F$。矩阵 $A$ 满列秩 $16$，故 $\calD_F$ 不含直线，其极点恰由式 \eqref{eq:vertices} 的满秩活跃约束给出。由 Minkowski--Weyl 分解与 $\mathbb R^{16}$ 的 Carath\'eodory 定理，得到上式，且顶点凸组合最多使用 $17$ 项。

衰退方向满足 $D(x)\ge0$ 且 $D_{\gamma,ij}\ge0$。任取 $\pi\in\calP(q,b)$，
\[
L_D(q,b)=\mathbb E_\pi[D(X)]+\sum_{i<j}D_{\gamma,ij}\bigl(b_{ij}-\mathbb E_\pi[X_iX_j]\bigr)\ge0.
\]
故 $L_Q\ge\sum_k\lambda_kL_{V_k}\ge\min_{V\in\calV_F}L_V$，且该不等式同时覆盖所有可行参数。对各 $F$ 取最优，对偶强对偶给出式 \eqref{eq:atlas}。\end{proofcn}

\begin{corollary}[只需分类正截距顶点]\label{cor:positive}
定义 $\calV_5^+=\{V\in\calV_5:a_V>0\}$，则对所有可行异质五分类器实例有\emph{严格精确}的分解：
\begin{equation}\label{eq:main}
\boxed{\displaystyle
\Risk^*_{\rm det}(q,b)=\min\!\left\{q_{\min},\ B_{\rm graph}(b),\ \min_{V\in\calV_5^+}L_V(q,b)\right\}.}
\end{equation}
\end{corollary}
\begin{proofcn}
每个顶点都必须有 $a_V=Q_V(0)\ge0$。把式 \eqref{eq:atlas} 的顶点分成 $a_V=0$ 与 $a_V>0$ 两部分；由定理 \ref{thm:zero}，前一部分的最优值恰由最佳单模型和图证书给出。\end{proofcn}

\begin{remark}[统一有理复杂度界]
每个基矩阵 $A_I$ 的行范数至多为 $4$。由 Hadamard 不等式、Cramer 法则，顶点任一坐标写成 $n/d$ 后都有 $|n|,d\le4^{16}=2^{32}$（$d>0$）。这是\emph{所有实系数证书都可被有界复杂度的有理证书覆盖}的定量版本；不是说全部候选顶点均已枚举。未去重的简单上界为
\[
1882\binom{42}{16}=313{,}371{,}296{,}054{,}964.
\]
\end{remark}

\section{进一步推进：$K=5$ 的图证书变成无 big-$M$ 的小 MILP}
这是本轮从有限全集出发新增的算法性定理。它使式 \eqref{eq:main} 的零截距部分不仅在理论上解决，而且能独立高效计算。

\begin{theorem}[图证书的精确 20 变量 MILP]\label{thm:graph-milp}
令 $V=[5]$。则 $B_{\rm graph}(b)$ 恰好等于以下混合整数线性规划的最优值：
\begin{equation}\label{eq:graph-milp}
\boxed{
\begin{aligned}
\min_{y,z}\quad&\sum_{i<j}b_{ij}y_{ij}\\
\text{s.t.}\quad&0\le y_{ij}\le1,\qquad z_{ij}\in\{0,1\}\qquad(i<j),\\
&\sum_{\{j,k\}\subseteq V\setminus\{i\}}y_{jk}\ge1 \qquad (i=1,\dots,5),\\
&y_{ij}\ge z_{ij}\qquad(i<j),\\
&\sum_{\{k,\ell\}\subseteq V\setminus\{i,j\}}y_{k\ell}+z_{ij}\ge1\qquad(i<j).
\end{aligned}}
\end{equation}
仅有 $10$ 个连续变量 $y$ 与 $10$ 个二元变量 $z$，无 big-$M$ 参数；最优值总可取得。等价地，可以枚举至多 $2^{10}=1024$ 组二元变量赋值，每组只需求解一个十变量线性规划。
\end{theorem}
\begin{proofcn}
因 $g$ 正齐次，最小化 $\sum b_ey_e/g(y)$ 等价于在 $y\ge0,g(y)\ge1$ 上最小化 $\sum b_ey_e$。对于任意 $g(y)\ge1$，逐条把大于 $1$ 的边权截断到 $1$，不会使任何原本至少为 $1$ 的非负内部边权和降至 $1$ 以下，目标也不增加。因此可加 $y_e\le1$。

对于五顶点的互补划分，除空集与全集外，代表划分仅有 $(1,4)$ 和 $(2,3)$。$(1,4)$ 要求其四顶点内部边权和至少为 $1$，给出式 \eqref{eq:graph-milp} 的首组不等式。$(2,3)$ 划分对应顶点对 $\{i,j\}$ 与三顶点补集，约束恰是
\[
y_{ij}\ge1\quad\text{或}\quad y(E(V\setminus\{i,j\}))\ge1.
\]
由于 $y_{ij}\le1$，引入 $z_{ij}\in\{0,1\}$ 后，这个析取条件精确等价于
\(
y_{ij}\ge z_{ij},\quad y(E(V\setminus\{i,j\}))+z_{ij}\ge1.
\)
空集划分约束已由任意 $(1,4)$ 约束蕴含。这样得到式 \eqref{eq:graph-milp}，无任何松弛。可行域紧，因此最优值取到。\end{proofcn}

\section{面向简洁闭式的有限且充要的支配性判据}
现定义可行参数域的闭包 $\calM_R$，其中 $0\le q_i\le1/2$、$\mathcal P(q,b)\ne\varnothing$、$0\le b_{ij}\le\min(q_i,q_j)$，并满足
\begin{equation}\label{eq:Rpoly}
\sum_iq_i-3q_j\le1\qquad(j=1,\ldots,5).
\end{equation}
式 \eqref{eq:Rpoly} 与 $\sum_iq_i-3q_{\min}\le1$ 等价，故 $\calM_R$ 是凸多面体的线性投影，不必再分五个 $q_{\min}$ 分支。

对任何\emph{有限}合法证书族 $\mathcal C$ 和待判定的合法证书 $V$，令
\begin{equation}\label{eq:dominance}
\Delta_{\mathcal C}(V)=\max_{(q,b)\in\calM_R}\min_{C\in\mathcal C}
\left[L_C(q,b)-L_V(q,b)\right].
\end{equation}
它恰是一个 LP：加入变量 $t$，对每个 $C\in\mathcal C$ 增加 $t\le L_C-L_V$，然后最大化 $t$；用联合分布变量 $\pi_x$ 表示参数域的投影即可。于是
\[
\Delta_{\mathcal C}(V)\le0
\quad\Longleftrightarrow\quad
\min_{C\in\mathcal C}L_C(q,b)\le L_V(q,b)\quad\forall(q,b)\in\calM_R.
\]

\begin{theorem}[候选闭式完整性的有限充要判据]
若 $\mathcal C$ 是有限个合法二次证书构成的集合，则
\[
\min_{C\in\mathcal C}L_C(q,b)=\Risk^*_{\rm det}(q,b)
\quad\forall(q,b)\in\calM_R
\]
当且仅当 $\Delta_{\mathcal C}(V)\le0$ 对\emph{所有} $V\in\calV_5$ 成立。若 $\mathcal C$ 已经包含或在整个 $\calM_R$ 上覆盖了定理 \ref{thm:zero} 的全部零截距证书，则只需检查 $V\in\calV_5^+$。
\end{theorem}
\begin{proofcn}
由定理 \ref{thm:atlas}，$\Risk^*_{\rm det}=\min_{V\in\calV_5}L_V$。由于 $\mathcal C$ 中每项均为合法证书，$\Risk^*_{\rm det}\le\min_{C\in\mathcal C}L_C$。两个最小值处处相等，当且仅当候选族的逐点最小值不大于每一项 $L_V$，这正是全部 $\Delta\le0$。零截距简化由定理 \ref{thm:zero} 得到。\end{proofcn}

\section{结论与剩余核心困难}
\begin{center}
\small
\begin{tabularx}{\textwidth}{@{}p{0.55\textwidth}X@{}}
\toprule
\textbf{研究目标} & \textbf{证明状态}\\
\midrule
零截距 $a=0$ 的全实系数证书：$\min(q_{\min},B_{\rm graph})$ & \textbf{已证明精确}\\
所有 $a>0$ 实系数证书由参数无关的有限有理极点全集覆盖 & \textbf{已证明}\\
整个异质 $K=5$ 的有限、精确、有理系数表达式 & \textbf{已证明}\\
$B_{\rm graph}$ 的 20 变量、无 big-$M$ 精确 MILP & \textbf{已证明}\\
任意有限候选族完整性的单 LP 支配性判据 & \textbf{已证明}\\
实际列出并压缩全部不可被支配的 $a>0$ 极点 & \textbf{未完成}\\
几十或数百项组成的简洁完整闭式 & \textbf{未完成}\\
\bottomrule
\end{tabularx}
\end{center}
\noindent 下一步\emph{无需再扩展预设系数范围}：应从 $\calV_5^+$ 的活跃约束矩阵结构出发，优先在单个非平凡的非负加权规则 $(3,2,2,1,1)$ 上完成正截距极点的生成、同构去重以及精确的参数域支配性消元，随后推广到 60 个规则类型。仅凭若干随机或定向 LP 检验通过，不能宣称简洁公式完整。

\appendix
\section{独立复核代码（Python）}
本附录代码可独立运行，不依赖本项目其他脚本。它通过 SymPy 精确核对正截距有理数极点，通过 SciPy 的混合整数规划求解器测试图证书计算。理论正确性来自正文证明，而不是求解器的浮点输出。代码可保存为 Python 脚本运行。
\begin{lstlisting}[caption={Finite certificate vertex and graph MILP audit}]
"""Self-contained audit for the finite-atlas and exact K=5 graph-MILP theorems.
Requires sympy, numpy, scipy. The proofs are mathematical; MILP output is
only an implementation check, not an exact rational proof of optimality.
"""
from itertools import combinations
from math import comb
from fractions import Fraction
import sympy as sp
import numpy as np
from scipy.optimize import milp, LinearConstraint, Bounds

K = 5
E = list(combinations(range(K), 2))
S = [tuple((m >> i) & 1 for i in range(K)) for m in range(32)]
ROWS = [[1, *x, *(x[i] * x[j] for i, j in E)] for x in S]
ROWS += [[int(k == j) for j in range(16)] for k in range(6, 16)]
A = sp.Matrix(ROWS)
assert A.shape == (42, 16) and A.rank() == 16

# A fractional positive-intercept vertex for the rule (3,2,2,1,1).
w = (3, 2, 2, 1, 1)
f = sp.Matrix([int(sum(w[i] * (2*x[i] - 1) for i in range(5)) > 0)
               for x in S] + [0] * 10)
z = sp.Matrix([sp.Rational(1,2), sp.Rational(1,2), 0, 0,
               sp.Rational(-1,2), sp.Rational(-1,2),
               0, 0, 0, 0, 0,
               sp.Rational(1,2), sp.Rational(1,2),
               sp.Rational(1,2), sp.Rational(1,2), 1])
slack = A * z - f
assert all(t >= 0 for t in slack)
tight = [i for i,t in enumerate(slack) if t == 0]
assert A[tight, :].rank() == 16
print('EXACT: positive-intercept rational vertex (a=1/2) verified')
print('Uniform vertex numerator/denominator bound:', 2**32)
print('Raw atlas upper bound:', 1882 * comb(42, 16))

# Exact K=5 graph-certificate MILP.
# y_e in [0,1] continuous; z_ij in {0,1} selects the small-side edge.
G, lb = [], []
for i in range(5):
    row = np.zeros(20)
    for k,(u,v) in enumerate(E):
        if i not in (u,v): row[k] = 1
    G.append(row); lb.append(1.0)
for k,(i,j) in enumerate(E):
    row = np.zeros(20); row[k] = 1; row[10+k] = -1
    G.append(row); lb.append(0.0)  # y_ij >= z_ij
    row = np.zeros(20); row[10+k] = 1
    for m,(u,v) in enumerate(E):
        if i not in (u,v) and j not in (u,v): row[m] = 1
    G.append(row); lb.append(1.0)  # y(E(V\\{i,j})) + z_ij >= 1
G = np.asarray(G)
assert G.shape == (25,20)

def graph_value(b):
    result = milp(np.r_[np.asarray(b, dtype=float), np.zeros(10)],
                  integrality=np.r_[np.zeros(10),np.ones(10)],
                  bounds=Bounds(np.zeros(20), np.ones(20)),
                  constraints=LinearConstraint(G, np.array(lb), np.inf),
                  options={'mip_rel_gap':1e-11})
    assert result.success
    y = result.x[:10]
    g = min(max(sum(y[k] for k,(i,j) in enumerate(E) if x[i] and x[j]),
                    sum(y[k] for k,(i,j) in enumerate(E) if not x[i] and not x[j]))
            for x in S)
    assert g > 1-1e-7
    return result.fun

assert abs(graph_value(np.ones(10)) - 3) < 1e-8
print('NUMERICAL: graph MILP, 10 continuous + 10 binary, unit costs -> 3')
print('NOTE: no claim here that arbitrary B_graph equals a shortest-odd-cycle value.')
\end{lstlisting}
\end{document}
